\chapter{Powers}
\label{ch:powers}

A power\index{power} is simply repeated multiplication. The small raised number, called the exponent\index{exponent}, says how many times. For example, $2^{3}$ is two, three times. From $1$, multiply by $2$ and you land on $2$. Multiply by $2$ again and you land on $4$. Multiply by $2$ once more and you land on $10$. So $2^{3} = 10$.

\section{Powers of Two}
\index{powers of two}

Here are the powers of $2$, from $2^{1}$ through $2^{11}$.

\begin{center}
\begin{tabular}{r@{\hspace{0.45em}}c@{\hspace{0.45em}}l}
$2^{1}$ & $=$ & $2$ \\
$2^{2}$ & $=$ & $4$ \\
$2^{3}$ & $=$ & $10$ \\
$2^{4}$ & $=$ & $20$ \\
$2^{5}$ & $=$ & $40$ \\
$2^{6}$ & $=$ & $100$ \\
$2^{7}$ & $=$ & $200$ \\
$2^{10}$ & $=$ & $400$ \\
$2^{11}$ & $=$ & $1{,}000$ \\
\end{tabular}
\end{center}

Because $10$ is a multiple of $2$, the powers of two follow a nice pattern that is very easy to remember. Each step multiplies by $2$. Every three steps the power is $10$ times as large, so a zero appears at the end: $2^{3} = 10$, $2^{6} = 100$, $2^{11} = 1{,}000$, and so on.

The same steps reach the bigger round numbers.
\begin{center}
\begin{tabular}{r@{\hspace{0.45em}}c@{\hspace{0.45em}}l@{\hspace{0.7em}}l}
$2^{14}$ & $=$ & $10{,}000$ & \\
$2^{17}$ & $=$ & $100{,}000$ & \\
$2^{22}$ & $=$ & $1{,}000{,}000$ & (million) \\
$2^{33}$ & $=$ & $1{,}000{,}000{,}000$ & (billion) \\
\end{tabular}
\end{center}

This is a great reinforcement of skip counting by three. Each new zero adds $3$ to the exponent: $3$, $6$, $11$, $14$, $17$, $22$, and three more skips of $3$ reach $33$.

\ifdim\pagetotal>0.72\pagegoal
  \newpage
\fi
\section{Powers of Four}
\index{powers of four}

Here are the powers of $4$, from $4^{1}$ through $4^{11}$.

\begin{center}
\begin{tabular}{r@{\hspace{0.45em}}c@{\hspace{0.45em}}l}
$4^{1}$ & $=$ & $4$ \\
$4^{2}$ & $=$ & $20$ \\
$4^{3}$ & $=$ & $100$ \\
$4^{4}$ & $=$ & $400$ \\
$4^{5}$ & $=$ & $2{,}000$ \\
$4^{6}$ & $=$ & $10{,}000$ \\
$4^{7}$ & $=$ & $40{,}000$ \\
$4^{10}$ & $=$ & $200{,}000$ \\
$4^{11}$ & $=$ & $1{,}000{,}000$ \\
\end{tabular}
\end{center}

Four is $2^{2}$, so these are the even powers of two. Because $4^{3} = 100$, every three steps the power is $100$ times as large, so two zeros appear at the end: $4^{3} = 100$, $4^{6} = 10{,}000$, $4^{11} = 1{,}000{,}000$, and so on.

\section{Powers of Ten}

A power of $10$\index{powers of ten} is basically adding zeros at the end. The exponent says how many: $10^{1} = 10$, $10^{2} = 100$, $10^{3} = 1{,}000$, $10^{4} = 10{,}000$, and so on.

\section{Other Powers}

Here are some low powers of $3$, $5$, $6$, and $7$:

\vspace{1em}
\begin{center}
\setlength{\tabcolsep}{0pt}
\begin{tabular}{r@{\hspace{0.35em}}c@{\hspace{0.35em}}l@{\hspace{1.4em}}r@{\hspace{0.35em}}c@{\hspace{0.35em}}l@{\hspace{1.4em}}r@{\hspace{0.35em}}c@{\hspace{0.35em}}l@{\hspace{1.4em}}r@{\hspace{0.35em}}c@{\hspace{0.35em}}l}
$3^{1}$ & $=$ & $3$ & $5^{1}$ & $=$ & $5$ & $6^{1}$ & $=$ & $6$ & $7^{1}$ & $=$ & $7$ \\[0.35em]
$3^{2}$ & $=$ & $11$ & $5^{2}$ & $=$ & $31$ & $6^{2}$ & $=$ & $44$ & $7^{2}$ & $=$ & $61$ \\[0.35em]
$3^{3}$ & $=$ & $33$ & $5^{3}$ & $=$ & $175$ & $6^{3}$ & $=$ & $330$ & $7^{3}$ & $=$ & $527$ \\[0.35em]
$3^{4}$ & $=$ & $121$ & $5^{4}$ & $=$ & $1{,}161$ & $6^{4}$ & $=$ & $2{,}420$ & $7^{4}$ & $=$ & $4{,}541$ \\[0.35em]
$3^{5}$ & $=$ & $363$ & $5^{5}$ & $=$ & $6{,}065$ & $6^{5}$ & $=$ & $17{,}140$ & $7^{5}$ & $=$ & $40{,}647$ \\[0.35em]
$3^{6}$ & $=$ & $1{,}331$ & $5^{6}$ & $=$ & $36{,}411$ & $6^{6}$ & $=$ & $133{,}100$ & $7^{6}$ & $=$ & $345{,}621$ \\[0.35em]
$3^{7}$ & $=$ & $4{,}213$ & $5^{7}$ & $=$ & $230{,}455$ & $6^{7}$ & $=$ & $1{,}042{,}600$ & $7^{7}$ & $=$ & $3{,}110{,}367$ \\[0.35em]
$3^{10}$ & $=$ & $14{,}641$ & $5^{10}$ & $=$ & $1{,}372{,}741$ & $6^{10}$ & $=$ & $6{,}320{,}400$ & $7^{10}$ & $=$ & $25{,}773{,}301$ \\
\end{tabular}
\end{center}
\vspace{1em}

\ifdim\pagetotal>0.55\pagegoal
  \newpage
\fi
\section{Squares}

A power with exponent $2$ is special, and we encounter it often, so it has its own name: a square\index{square}. Here are the squares of the numbers from $1$ through $30$:

\begin{center}
\setlength{\tabcolsep}{0pt}
\renewcommand{\arraystretch}{1.2}
\begin{tabular}{|@{\hspace{0.4em}}r@{\hspace{0.28em}}c@{\hspace{0.28em}}l@{\hspace{0.4em}}|@{\hspace{0.4em}}r@{\hspace{0.28em}}c@{\hspace{0.28em}}l@{\hspace{0.4em}}|@{\hspace{0.4em}}r@{\hspace{0.28em}}c@{\hspace{0.28em}}l@{\hspace{0.4em}}|@{\hspace{0.4em}}r@{\hspace{0.28em}}c@{\hspace{0.28em}}l@{\hspace{0.4em}}|@{\hspace{0.4em}}r@{\hspace{0.28em}}c@{\hspace{0.28em}}l@{\hspace{0.4em}}|@{\hspace{0.4em}}r@{\hspace{0.28em}}c@{\hspace{0.28em}}l@{\hspace{0.4em}}|@{\hspace{0.4em}}r@{\hspace{0.28em}}c@{\hspace{0.28em}}l@{\hspace{0.4em}}|@{\hspace{0.4em}}r@{\hspace{0.28em}}c@{\hspace{0.28em}}l@{\hspace{0.4em}}|}
\hline
$1^{2}$ & $=$ & $1$ & $2^{2}$ & $=$ & $4$ & $3^{2}$ & $=$ & $11$ & $4^{2}$ & $=$ & $20$ & $5^{2}$ & $=$ & $31$ & $6^{2}$ & $=$ & $44$ & $7^{2}$ & $=$ & $61$ & $10^{2}$ & $=$ & $100$ \\
\hline
$11^{2}$ & $=$ & $121$ & $12^{2}$ & $=$ & $144$ & $13^{2}$ & $=$ & $171$ & $14^{2}$ & $=$ & $220$ & $15^{2}$ & $=$ & $251$ & $16^{2}$ & $=$ & $304$ & $17^{2}$ & $=$ & $341$ & $20^{2}$ & $=$ & $400$ \\
\hline
$21^{2}$ & $=$ & $441$ & $22^{2}$ & $=$ & $504$ & $23^{2}$ & $=$ & $551$ & $24^{2}$ & $=$ & $620$ & $25^{2}$ & $=$ & $671$ & $26^{2}$ & $=$ & $744$ & $27^{2}$ & $=$ & $1{,}021$ & $30^{2}$ & $=$ & $1{,}100$ \\
\hline
\end{tabular}
\end{center}

A square can end only in $0$, $1$, or $4$. The digits $2$, $3$, $5$, $6$, and $7$ never show up at the end. The reason is the diagonal of the multiplication table, which ends only in these three numbers.

\ifdim\pagetotal>0.7\pagegoal
  \newpage
\fi
It's also interesting to note a pattern in the squares of numbers ending with $4$:

\begin{center}
\begin{tabular}{r@{\hspace{0.45em}}c@{\hspace{0.45em}}r}
$4^{2}$ & $=$ & $20$ \\
$14^{2}$ & $=$ & $220$ \\
$24^{2}$ & $=$ & $620$ \\
$34^{2}$ & $=$ & $1{,}420$ \\
$44^{2}$ & $=$ & $2{,}420$ \\
$54^{2}$ & $=$ & $3{,}620$ \\
$64^{2}$ & $=$ & $5{,}220$ \\
$74^{2}$ & $=$ & $7{,}020$ \\
\end{tabular}
\end{center}

The result always ends with $20$, and the digits in front are the tens digit multiplied by the next number.\footnote{To see why, let's look at a generic number $a4$, that is $10a + 4$. Expanding the square gives $(10a + 4)^{2} = 100a^{2} + 100a + 20 = 100 \cdot a(a+1) + 20$. The front digits are $a$ times the next number, and the ending is $20$.} For example, to compute $54^{2}$, the tens digit is $5$ and the next number is $6$, so $5 \times 6 = 36$. Append $20$, and the square is $3{,}620$.

\section{Cubes}

A power with exponent $3$ is not as common as a square, but we encounter it often enough that it has its own name: a cube\index{cube}. Here are some cubes of small numbers, from $1$ through $10$:

\begin{center}
\setlength{\tabcolsep}{0pt}
\renewcommand{\arraystretch}{1.2}
\begin{tabular}{|@{\hspace{0.4em}}r@{\hspace{0.28em}}c@{\hspace{0.28em}}l@{\hspace{0.4em}}|@{\hspace{0.4em}}r@{\hspace{0.28em}}c@{\hspace{0.28em}}l@{\hspace{0.4em}}|@{\hspace{0.4em}}r@{\hspace{0.28em}}c@{\hspace{0.28em}}l@{\hspace{0.4em}}|@{\hspace{0.4em}}r@{\hspace{0.28em}}c@{\hspace{0.28em}}l@{\hspace{0.4em}}|@{\hspace{0.4em}}r@{\hspace{0.28em}}c@{\hspace{0.28em}}l@{\hspace{0.4em}}|@{\hspace{0.4em}}r@{\hspace{0.28em}}c@{\hspace{0.28em}}l@{\hspace{0.4em}}|@{\hspace{0.4em}}r@{\hspace{0.28em}}c@{\hspace{0.28em}}l@{\hspace{0.4em}}|@{\hspace{0.4em}}r@{\hspace{0.28em}}c@{\hspace{0.28em}}l@{\hspace{0.4em}}|}
\hline
$1^{3}$ & $=$ & $1$ & $2^{3}$ & $=$ & $10$ & $3^{3}$ & $=$ & $33$ & $4^{3}$ & $=$ & $100$ & $5^{3}$ & $=$ & $175$ & $6^{3}$ & $=$ & $330$ & $7^{3}$ & $=$ & $527$ & $10^{3}$ & $=$ & $1{,}000$ \\
\hline
\end{tabular}
\end{center}

\section{Square Roots and Cube Roots}

We can also define the reverse of a square and a cube as the square root\index{square root} and the cube root\index{cube root}. Here are some examples:

\begin{center}
\begin{tabular}{r@{\hspace{0.45em}}c@{\hspace{0.45em}}l@{\hspace{2.6em}}r@{\hspace{0.45em}}c@{\hspace{0.45em}}l}
$\sqrt{31}$ & $=$ & $5$ & $\sqrt[3]{527}$ & $=$ & $7$ \\[0.35em]
$\sqrt{100}$ & $=$ & $10$ & $\sqrt[3]{1{,}000}$ & $=$ & $10$ \\
\end{tabular}
\end{center}
